% =============================================================================
%  APÊNDICE B: GLOSSÁRIO DE TERMOS TÉCNICOS
% =============================================================================
\chapter{Glossário de Termos Técnicos}
\label{apend:glossario}

Este glossário reúne os principais termos técnicos utilizados ao longo do
livro. Cada verbete apresenta o termo em negrito, sua definição concisa e
a seção ou capítulo do livro onde é discutido em profundidade.

\begin{description}

\item[Agente Cognitivo] Programa de software com capacidade de percepção,
raciocínio e ação autônoma em um ambiente. No \opencodeeco{}, 128 agentes
especializados formam o núcleo operacional. (Capítulo~2)

\item[AlphaFold] Sistema de IA do Google DeepMind para predição de
estrutura tridimensional de proteínas, integrado como skill científica
no \opencodeeco{}. (Seção~2.3)

\item[Arquitetura Três Camadas] Estrutura fundamental do \opencodeeco{}
organizada em três camadas: MCP, Skill e Agent. (Capítulo~3)

\item[Auto-Consciência Artificial] Capacidade de um sistema de IA de
construir e manter um modelo de si mesmo, classificado nos níveis
N0 (inconsciente) a N3 (auto-consciente pleno) no \opencodeeco{}.
(Capítulo~4)

\item[Auto-Evolução] Processo pelo qual o ecossistema identifica lacunas
em suas próprias capacidades e as preenche autonomamente, guiado pelo
Scanner Pipeline. (Capítulo~4)

\item[Barramento de Eventos] Mecanismo de comunicação assíncrona entre
componentes do ecossistema, baseado no padrão publish-subscribe.
(Seção~3.4)

\item[Behavioral Gate] Componente do Trust Engine que classifica ações
como seguras, moderadas, arriscadas ou bloqueadas antes da execução.
(Seção~5.3)

\item[BettaFish] Conjunto de 11 ferramentas de análise acadêmica
(OASIS, Forum, Config, Graph, Report, Nash, Stats, Qualis,
Sensitivity, IMRAD, Debate) integradas ao pipeline P14--P18 do
\opencodeeco{}. (Seção~3.8)

\item[Capability Composer] Scanner que decompõe capacidades complexas em
insumos cognitivos atômicos, aplicando desconto por compartilhamento.
(Seção~4.6)

\item[Ciclo PLAN--ACT--REFLECT--EXTRACT--EVOLVE] Ciclo evolutivo do
Manus Evolve que guia a descoberta e instalação de novas skills.
(Seção~3.6)

\item[Composição Unitária do Conhecimento] Metodologia de decomposição
de capacidades em 6 tipos de insumos cognitivos atômicos (SPEC-033).
(Seção~4.6)

\item[CORA-Eval] Benchmark de 150 tarefas em 10 dimensões e 4 níveis
para avaliar ecossistemas cognitivos. (Capítulo~7)

\item[Dialectical Engine] Motor dialético que implementa a tríade
tese--antítese--síntese para resolução de conflitos entre agentes.
(Seção~4.9)

\item[Entropia de Shannon] Medida de incerteza ou conteúdo informacional
de uma fonte, medida em bits. (Seção~1.7)

\item[Evolutionary Trajectories Scanner] Scanner que projeta e avalia
diferentes rotas evolutivas para o ecossistema. (Seção~4.4)

\item[Gradiente Descendente] Algoritmo de otimização iterativa que
minimiza funções seguindo a direção oposta ao gradiente. (Seção~1.4)

\item[Governança Cooperativa] Framework baseado nos 8 Princípios de
Ostrom (DP1--DP8) para governança descentralizada de agentes.
(Seção~4.9)

\item[Injeção de Dependência] Padrão de design que fornece dependências
a um componente externamente, promovendo baixo acoplamento.
(Seção~3.4)

\item[Insumo Cognitivo] Unidade atômica de conhecimento nos 6 tipos:
dado, informação, conceito, procedimento, heurística e axioma.
(Seção~4.6)

\item[LLM] \emph{Large Language Model} -- modelo de linguagem de grande
escala, como GPT-4 e DeepSeek, utilizado como motor de raciocínio
central no \opencodeeco{}. (Seção~2.4)

\item[Lógica Proposicional] Sistema formal que estuda proposições e
conectivos lógicos (E, OU, NÃO, SE\;...\;ENTÃO). (Seção~1.1)

\item[MASWOS] \emph{Multi-Agent System for Writing and Orienting
Scientific} -- sistema multiagente para produção acadêmica com 49
agentes especializados. (Capítulo~8)

\item[MCSP] \emph{Minimum Capability Set Problem} -- problema de
seleção do conjunto mínimo de capacidades para cobrir funcionalidades
requeridas (SPEC-032). (Seção~4.5)

\item[MCP] \emph{Model Context Protocol} -- protocolo padrão para
comunicação entre modelos de IA e ferramentas externas. O
\opencodeeco{} integra 46 servidores MCP. (Capítulo~3)

\item[Memória Atkinson-Shiffrin] Modelo de memória em três estágios:
sensorial, curto prazo e longo prazo, adotado pelo
\texttt{NaturalForgetting}. (Seção~5.4)

\item[Metacognição] Camada de auto-observação do ecossistema composta
por MetacognitiveMonitor, DialecticalEngine, CooperativeGovernance e
SelfModel (SPEC-036). (Seção~4.8)

\item[MetacognitiveMonitor] Componente que supervisiona continuamente
o estado interno do ecossistema, detectando anomalias e lacunas.
(Seção~4.8)

\item[MiroFish] Conjunto de ferramentas de modelagem e simulação
multiagente integradas ao pipeline P14--P18. (Seção~3.8)

\item[Modelo N3] Terceiro nível de auto-consciência artificial, onde o
sistema mantém um modelo completo de si mesmo e de seus limites.
(Seção~4.8)

\item[N3.5] Extensão do N3 que adiciona o Behavioral Gate preventivo,
atingido no ciclo R23 do \opencodeeco{}. (Capítulo~5)

\item[NaturalForgetting] Mecanismo de esquecimento natural baseado no
modelo Atkinson-Shiffrin que remove gradualmente informações obsoletas.
(Seção~5.4)

\item[Nexus] Camada de orquestração multiagente com 488 arquivos,
120+ barreiras de sincronização e 212+ tipos de raciocínio.
(Seção~3.7)

\item[Noological Scanner] Scanner que identifica lacunas de capacidade
no ecossistema perguntando ``o que não existe?'' (SPEC-028).
(Seção~4.2)

\item[OpenCode Ecosystem] Plataforma de engenharia de software com
agentes inteligentes, integrando 46 MCPs, 227 skills, 128 agentes,
15 plug-ins e 14 comandos. (Capítulo~3)

\item[OutcomeTracker] Componente que registra o resultado de cada ação
executada por um agente, alimentando o TrustScorer. (Seção~5.5)

\item[PhD Auditor] Sistema de auditoria acadêmica com 6 módulos:
NashSolver, StatisticalRigor, QualisA1Auditor, SensitivityAnalyzer,
IMRADFormatter e CrossValidation. (Seção~3.8)

\item[Potentiality Scanner] Scanner que identifica capacidades
emergentes no ecossistema (SPEC-043). (Seção~4.7)

\item[Qualis A1] Classificação máxima no sistema Qualis da CAPES para
periódicos científicos. (Capítulo~8)

\item[Raciocínio Dialético] Método de argumentação baseado na tríade
tese--antítese--síntese, implementado no DialecticalEngine.
(Seção~4.9)

\item[Scanner Pipeline] Conjunto de 7 scanners epistemológicos
encadeados: Noological, Teleological Reverse, Evolutionary Trajectories,
Refinement, MCSP, Capability Composer e Potentiality. (Capítulo~4)

\item[Scanner Refinement] Scanner que refina e valida as capacidades
propostas pelos scanners anteriores (SPEC-031). (Seção~4.5)

\item[SDD] \emph{Spec-Driven Development} -- metodologia em que a
especificação formal precede a implementação. (Seção~3.3)

\item[SEEKER] Sistema de pesquisa com 10 agentes inteligentes e motor
de árvore de argumentos para fundamentação acadêmica.
(Capítulo~8)

\item[SelfModel] Modelo que o ecossistema mantém de si mesmo,
evoluindo dos níveis N0 a N3. (Seção~4.8)

\item[Skill] Habilidade especializada no \opencodeeco{}, totalizando
227 skills em 13 categorias. (Seção~3.6)

\item[Structural Noise Scanner] Scanner que aplica compressão estrutural
com preservação funcional (SPEC-037). (Capítulo~5)

\item[TDD] \emph{Test-Driven Development} -- metodologia em que os
testes são escritos antes do código de produção. (Seção~3.3)

\item[Teleological Reverse Scanner] Scanner que projeta o estado futuro
desejado e trabalha reversamente para identificar capacidades
necessárias (SPEC-029). (Seção~4.3)

\item[Token Economy] Sistema de incentivos econômicos baseado em tokens
para recompensar contribuições ao ecossistema (SPEC-022 a SPEC-024).
(Capítulo~6)

\item[Transformer] Arquitetura de rede neural baseada em mecanismos de
atenção, fundamento dos LLMs modernos. (Seção~2.4)

\item[Trust Engine] Sistema integrado de pontuação de confiança, barreiras
comportamentais, esquecimento natural e governança (SPEC-038).
(Capítulo~5)

\item[TrustScorer] Componente do Trust Engine que calcula a pontuação
de confiança de cada agente usando blend 70/30 entre ações diretas e
feedback da rede. (Seção~5.2)

\item[Z3] Provador de teoremas e solver SMT da Microsoft, integrado
como motor de raciocínio formal no \opencodeeco{}. (Seção~2.7)

\item[SMT] \emph{Satisfiability Modulo Theories} -- extensão do SAT
para teorias de primeira ordem. (Seção~2.7)

\item[SymPy] Biblioteca Python para matemática simbólica, integrada
como motor de raciocínio simbólico. (Seção~2.7)

\item[MiniKanren] Linguagem de programação lógica relacional, integrada
como motor de raciocínio lógico. (Seção~2.7)

\item[Critical Reasoning] Motor de análise de falácias lógicas e vieses
cognitivos, com 15 tipos de falácias. (Seção~2.7)

\item[Manus Evolve] Motor de evolução autônoma que gerencia o ciclo
PLAN--ACT--REFLECT--EXTRACT--EVOLVE. (Seção~3.6)

\item[Shadow Mode] Modo de operação monitorada no Trust Engine, ativado
quando a confiança de um agente cai abaixo do limiar. (Seção~5.2)

\item[Staking] Mecanismo de bloqueio de tokens por 7 dias para
participar da governança do ecossistema. (Capítulo~6)

\item[Slashing] Penalidade de redução de stake por comportamento
malicioso ou negligente. (Capítulo~6)

\item[Fee Market] Mercado dinâmico de taxas para uso de recursos do
ecossistema, com preços ajustados por oferta e demanda.
(Capítulo~6)

\item[Ledger] Registro imutável de transações no sistema de Token
Economy, implementado como \texttt{frozen dataclass} em Python.
(Capítulo~6)

\item[Audit Trail] Trilha de auditoria com hash SHA-256 para todas as
transações do sistema. (Capítulo~6)

\item[Plugin] Módulo extensível que adiciona funcionalidades ao
ecossistema; são 15 plug-ins (10 npm, 2 locais .ts, 3 bridge).
(Seção~3.5)

\item[Quantum Nexus] Módulo de computação quântica com 146 arquivos,
incluindo QML para HAM10000 (89,52\%) e 50 qubits MPS.
(Seção~7.4)

\item[Granger Causality] Teste estatístico para determinar se uma série
temporal ajuda a predizer outra, usado no N3. (Seção~5.6)

\item[Bayesian Inference] Método de inferência estatística baseado no
Teorema de Bayes, usado no diagnóstico do Trust Engine.
(Seção~5.6)

\item[Ostrom Principles] Conjunto de 8 princípios de governança de
recursos comuns (DP1--DP8), aplicados à governança de agentes.
(Seção~4.9)

\item[ADR] \emph{Architecture Decision Record} -- registro de decisões
arquiteturais; 10 ADRs documentadas no \opencodeeco{}.
(Seção~3.9)

\item[Evolução Autônoma] Capacidade do ecossistema de gerar novas skills
sem intervenção humana, através do Manus Evolve. (Seção~3.6)

\item[LSP] \emph{Language Server Protocol} -- protocolo para servidores
de linguagem; TypeScript LSP integrado. (Seção~3.4)

\item[Agente Autônomo] Agente com capacidade de operar sem supervisão
humana direta, dentro dos limites do Trust Engine. (Seção~2.8)

\item[Argument Tree] Estrutura de árvore que organiza argumentos e
contra-argumentos, usada pelo SEEKER. (Capítulo~8)

\item[TSAC] \emph{Text Style Anti-Cloning} -- técnica de detecção e
prevenção de textos similares a saídas de IA, com 87 palavras
proibidas. (Capítulo~8)

\item[Cross-Validation] Técnica de validação de modelos que particiona
os dados em k subconjuntos, usando k-1 para treino e 1 para teste.
(Capítulo~7)

\item[Teoria da Informação] Ramo da matemática que estuda a quantificação,
armazenamento e comunicação de informação. (Seção~1.7)

\item[Teoria dos Grafos] Ramo da matemática que estuda estruturas
compostas por vértices e arestas. (Seção~1.8)

\item[Complexidade Computacional] Ramo da ciência da computação que
classifica problemas segundo sua dificuldade inerente. (Seção~1.9)

\item[Lema] Proposição auxiliar utilizada como passo intermediário na
demonstração de um teorema. (Capítulo~1)

\item[Corolário] Proposição que decorre diretamente de um teorema já
demonstrado. (Capítulo~1)

\item[TikZ] Pacote LaTeX para criação de gráficos vetoriais e diagramas
programaticamente. (Usado em todos os capítulos)

\item[Hyperparameter] Parâmetro cujo valor é definido antes do
treinamento de um modelo de aprendizado de máquina. (Seção~2.3)

\item[Overfitting] Fenômeno em que um modelo se ajusta excessivamente
aos dados de treino, perdendo capacidade de generalização.
(Seção~2.2)

\item[Teste de Turing] Teste comportamental para avaliar a capacidade
de uma máquina de exibir comportamento inteligente indistinguível
de um humano. (Seção~2.1)

\end{description}
